\documentclass[11pt,a4paper]{article}
\usepackage[T1]{fontenc}
\usepackage[utf8]{inputenc}
\usepackage[main=italian,provide=*]{babel}
\usepackage{amsmath,amssymb}
\usepackage{geometry}
\geometry{margin=2.5cm}
\usepackage{booktabs}
\usepackage{seqsplit}
\usepackage{hyperref}
\hypersetup{colorlinks=true,linkcolor=blue,citecolor=blue,urlcolor=blue}

\title{Benchmark esatto (Fase 3) — trimero a catena aperta}
\author{Note di lavoro — Tesi triennale, Samuele Schiavi\\ Relatore: Prof.\ Alessandro Chiesa}
\date{}

\begin{document}
\maketitle

\section*{Introduzione}

Questo documento raccoglie i moduli di benchmark esatto per il trimero di spin-1/2
($N=3$) in topologia a catena aperta $1$–$2$–$3$ (legami fisici $(1,2)$ e $(2,3)$,
nessun legame $(3,1)$). Sono i moduli su cui tutto il resto del lavoro (VQE, e in
futuro Trotter e correlazioni) si appoggia senza mai modificarli — stesso ruolo di
\texttt{dimer\_exact.py} per il dimero e di \texttt{trimer\_ring\_exact.py} per
l'anello.

Per ciascun file: cosa fa, con quali verifiche è stato controllato in questa sessione,
e cosa se ne può concludere — per i moduli \texttt{.py}, anche le funzioni interne e
come si usa il file.

\textbf{Nota.} \texttt{verifica\_catena\_preliminare.py} (ora rinominato
\texttt{verifica\_catena.py}) è stato spostato nel documento dedicato
\texttt{trimero\_catena\_verifiche.pdf}, insieme al suo gemello
\texttt{verifica\_stati\_catena.py} — non è più descritto qui.

\section*{\texttt{trimer\_chain\_exact.py}}

Analogo di \texttt{trimer\_ring\_exact.py} per la topologia a catena. Simmetria di
riflessione $1\leftrightarrow3$ (sulla coppia \emph{non} legata) al posto della
simmetria di scambio $1\leftrightarrow2$ dell'anello (sulla coppia legata) — stesso
metodo di Kambe, applicato a una coppia di classificazione diversa. Per $J>0$ il
fondamentale a basso campo è \textbf{sempre} il blocco $B$ (mai serve scegliere fra
$B$ e $C$ come nell'anello, dato che $E_B=-4J<E_C=0$ indipendentemente da $J$): una
semplificazione reale rispetto all'anello, non solo una topologia più povera.

\textbf{Verifiche.} Otto self-test base più cinque self-test DM. \textbf{Correzione
rispetto alla versione precedente di questo documento}: qui si affermava erroneamente
che il self-test di coerenza fra \texttt{exact\_sweep\_dm} e
\texttt{ground\_state\_projector\_dm} fosse già presente — non lo era, le due funzioni
non esistevano affatto nel modulo a inizio Fase 5 (il file conteneva solo
\texttt{dm\_term} e \texttt{trimer\_hamiltonian\_dm}, esplicitamente etichettate
``struttura disponibile, non usata''). Le due funzioni, più \texttt{dm\_min\_gap}
(mirror di quella dell'anello, con la stessa finestra $[b_c-w,b_c+w]$ per evitare sia
il gap a $b_c$ fisso sia la degenerazione di Kramers a $b=0$ — si veda
\texttt{analisi\_dm\_trimero\_anello.tex}, sez.\ Kramers), sono state aggiunte in
questa fase. I due self-test aggiuntivi (coerenza sweep/projector; linearità
$g_\text{min}(D)=2\sqrt6\,D$) sono stati eseguiti per la prima volta ora, non
``rieseguiti'': tutti e tredici i self-test sono superati a precisione macchina.

\textbf{Conclusione.} Stesso livello di maturità e affidabilità del modulo
dell'anello.

\textbf{Funzioni interne.}
\begin{itemize}
\item \texttt{trimer\_hamiltonian(J,\ b)} — l'Hamiltoniana completa come
\texttt{SparsePauliOp} (un solo parametro $J$, nessun $J'$).
\item \texttt{magnetization\_operator()} — $M_z=(Z_1+Z_2+Z_3)/2$.
\item \texttt{S13\_squared\_operator()} — il Casimir della coppia di
classificazione \emph{non legata} $(1,3)$ — l'analogo di
\texttt{S12\_squared\_operator()} dell'anello, ma sulla coppia diversa.
\item \texttt{S\_total\_squared\_operator()} — il Casimir totale $S^2$.
\item \texttt{analytic\_eigenvalues(J,\ b)} — le otto energie esatte
$E(S_{13},S,M)=E_\text{blocco}+2bM$, ordinate.
\item \texttt{critical\_field(J)} — il campo critico $b_c=3J$ dell'incrocio del
fondamentale (esiste solo per $J>0$).
\item \texttt{kambe\_states()} — dizionario \texttt{\{(blocco, M): vettore a 8
componenti\}} con gli otto stati espliciti.
\item \texttt{kambe\_energy(J,\ b,\ block,\ M)} — l'energia analitica del singolo
stato \texttt{(block, M)}.
\item \texttt{exact\_sweep(b\_values,\ J)} — diagonalizza $H$ su una griglia di
campo $b$.
\item \texttt{ground\_state\_projector(J,\ b,\ tol)} — il proiettore sul
multipletto fondamentale, robusto a eventuale degenerazione.
\item \texttt{dm\_term(i,\ j,\ D)} — il termine DM $D(X_iZ_j-Z_iX_j)$, stessa
forma già validata per dimero e anello.
\item \texttt{trimer\_hamiltonian\_dm(J,\ b,\ D)} — l'Hamiltoniana con il termine
DM simmetrico sotto $P_{13}$ ($D_{12}=-D_{23}=D$ — unica opzione simmetrica, per
costruzione, a differenza dell'anello dove esistono due opzioni A/B da confrontare).
\textbf{Attenzione}: l'analogia strutturale con l'Opzione A dell'anello (entrambe
conservano $S_{13}^2$/$S_{12}^2$ esattamente) \emph{non} implica lo stesso esito —
l'Opzione A dell'anello non apre mai il gap (incrocio $C\to A$, settori $S_{12}$
diversi), mentre qui il gap si apre (incrocio $B\to A$, \emph{stesso} settore
$S_{13}=1$): si veda \texttt{analisi\_dm\_trimero\_catena.tex}.
\item \texttt{exact\_sweep\_dm(b\_values,\ J,\ D)} — come \texttt{exact\_sweep} ma con
$H_\text{DM}$ acceso; media sul sottospazio degenere per $\langle M_z\rangle$ e gli
altri valori di aspettazione (necessaria: a $b=0$ il fondamentale è un doppietto di
Kramers per \emph{qualunque} $D$, si veda oltre).
\item \texttt{ground\_state\_projector\_dm(J,\ b,\ D,\ tol)} — come
\texttt{ground\_state\_projector} ma con $H_\text{DM}$ acceso.
\item {\sloppy\texttt{dm\_min\_gap(J,\ D,\ search\_half\_width)} — il gap vero
$g_\text{min}(D)=\min_{b\in[b_c-w,b_c+w]}[E_1(b;D)-E_0(b;D)]$ (formula ristretta a un
intorno di $b_c$, non a tutto l'asse $b$), per evitare sia l'artefatto del gap a
$b_c$ fisso sia la degenerazione di Kramers a $b=0$: a $b=0$ e per qualunque $D$ si
avrebbe $g_\text{min}=0$, mascherando l'intero effetto — si veda
\texttt{analisi\_dm\_trimero\_anello.tex}, sez.\ ``Una seconda trappola''.\par}
\item \texttt{\_P13\_matrix()} — la matrice $8\times8$ dello scambio dei qubit 1 e
3 (registro fisico, Famiglia 1).
\item \texttt{\_self\_test()} — gli otto self-test base.
\item \texttt{\_self\_test\_dm()} — cinque self-test: simmetria $P_{13}$ del termine
DM e sua rottura nel caso non simmetrico; conservazione di $S_{13}^2$; $D=0$
ripristina l'Hamiltoniana base; coerenza fra \texttt{exact\_sweep\_dm} e
\texttt{ground\_state\_projector\_dm}; linearità di $g_\text{min}(D)$ con pendenza
$2\sqrt6\approx4.899$.
\end{itemize}

\textbf{Come si usa.} Come libreria: \texttt{from trimer\_chain\_exact import
trimer\_hamiltonian, exact\_sweep, \ldots} — dipendenza comune a tutti i moduli VQE
della catena. Come script standalone: \texttt{python3 trimer\_chain\_exact.py}
esegue entrambi i blocchi di self-test e rigenera la figura
\texttt{trimer\_chain\_exact.png}.

\section*{\texttt{trimer\_chain\_exact.png}}

Figura generata dal modulo precedente, analoga a \texttt{trimer\_ring\_exact.png}
per l'anello: a sinistra lo spettro completo con il fondamentale colorato per blocco
(blu $=B$ sotto $b_c$, arancio $=A$ sopra); a destra la magnetizzazione del
fondamentale, con la stessa annotazione già usata per l'anello che distingue il punto
convenzionale $b=0$ (fondamentale degenere, $\langle M_z\rangle=0$ solo per
convenzione di media sul sottospazio, non un limite fisico) dall'incrocio vero a
$b_c=3J$ (salto discontinuo reale, da $-0.5$ a $-1.5$).

\textbf{Considerazioni.} File mancante nella prima stesura di questo documento,
segnalato dall'utente — colmato in questa revisione, come già fatto per
\texttt{animazione\_trimero\_catena\_aperta.html}.

\section*{\texttt{teoria\_trimero\_catena\_aperta.pdf}}

Teoria esatta della catena uniforme: simmetria $P_{13}$, Casimir $S_{13}^2$, spettro
chiuso a tre multipletti, campo critico $b_c=3J$, prova che gli incroci sono veri (non
anticrossing) — stessa profondità dimostrativa del documento equivalente per l'anello,
con la coppia di classificazione cambiata da $(1,2)$ a $(1,3)$.

\section*{\texttt{derivazione\_stati\_kambe\_trimero\_catena.pdf}}

Derivazione esplicita degli otto autostati, con una complicazione in più rispetto
all'anello: la coppia di classificazione $(1,3)$ \textbf{non è adiacente} nel registro
fisico $|q_1q_2q_3\rangle$, quindi la derivazione richiede una gestione esplicita del
riordino che nell'anello (coppia $(1,2)$, adiacente) non serve. Stesso metodo di fondo,
applicazione tecnicamente più delicata.

\textbf{Conclusione.} Il blocco $B$ risultante non è banale (pesi
$1/\sqrt6,1/\sqrt6,-\sqrt{2/3}$ su $|100\rangle,|001\rangle,|010\rangle$) — a
differenza del blocco a basso campo scelto nel punto di lavoro standard dell'anello
(blocco $C$, singoletto banale sulla coppia base). Conseguenza diretta per i moduli
VQE: nella catena non esiste una derivazione a mano semplice ($X,H,\mathrm{CX},X$) per
lo stato iniziale a basso campo, serve \texttt{prepare\_state} fin da subito.

\section*{\texttt{animazione\_trimero\_catena\_aperta.html}}

Visualizzazione interattiva, analoga a \texttt{animazione\_trimero\_isoscele.html}
per l'anello: geometria della catena, spettro con il fondamentale colorato per
blocco. \textbf{Nessuna mappa di fase 2D} (a differenza dell'anello): la catena ha un
solo parametro di accoppiamento $J$, non una coppia $(J,J')$, quindi non c'è un piano
su cui tracciare regioni — la sola dipendenza rilevante è da $J$ e $b$. Utile per
l'intuizione qualitativa, non per i numeri precisi (quelli vengono dal modulo esatto).

\textbf{Considerazioni.} File mancante nella prima stesura di questo documento —
l'anello aveva già la sua voce equivalente, la catena no: colmato in questa
revisione.

\section*{Conclusione generale}

Il modulo esatto della catena condivide con quello dell'anello metodo (decomposizione
di Kambe), convenzione qubit (Famiglia 1, $\text{site}_i\to\text{qubit}(N-i)$) e
livello di rigore (self-test a precisione macchina, verifica indipendente su griglia
estesa). Le differenze reali, non cosmetiche, rispetto all'anello:

\begin{itemize}
\item \textbf{Simmetria di classificazione diversa}: la catena classifica sulla coppia
\emph{non legata} $(1,3)$, l'anello sulla coppia \emph{legata} $(1,2)$ — stesso metodo
algebrico, applicato a una coppia fisicamente diversa.
\item \textbf{Un parametro in meno}: la catena ha una sola costante di accoppiamento
($J$), l'anello due ($J,J'$) — la mappa dei segni 2D dell'anello collassa a una sola
dimensione nella catena.
\item \textbf{Scelta del fondamentale a basso campo più semplice}: per $J>0$ è sempre
il blocco $B$, senza bisogno del confronto $E_B$ vs $E_C$ che serve invece
nell'anello.
\item \textbf{Blocco a basso campo più complesso}: il blocco $B$ della catena non è
mai banale, a differenza del blocco a basso campo del punto di lavoro standard
dell'anello — differenza che si propaga direttamente nei moduli VQE (documento
gemello sul VQE).
\end{itemize}

Il modulo è pronto a supportare tutto il lavoro VQE già costruito e le fasi ancora da
affrontare (Trotter, correlazioni dinamiche).

\end{document}
